\color{unidarkgreen}{\subsection{Контроль синхронизма и напряжений} \label{func:ksn}}
\color{black}

\begin{enumerate}

\item
Функциональный блок контроля синхронизма и напряжений (КСН) предназначен для формирования разрешающих сигналов, соответствующих условиям включения, для АПВ, ПАВ и оперативного включения.

\item
В данном РЭ под энергообъектом Э1 принято обозначение линии, а под энергообъектом Э2 -- система шин. В дальнейшем, по тексту, могут применяться оба обозначения.

\item
Основные режимы включения и контролируемые для них условия приведены в таблице \ref{ksn:tbl1}. Режимы включения с контролем синхронизма (КС) и с улавливанием синхронизма (УС) вводятся соответствующими программными ключами \textbf{«КС»} и \textbf{«УС»}.

\small
\setlength{\extrarowheight}{0.05cm} % высота ячеек в таблице
\begin{longtable}
{|>{\centering\arraybackslash}m{1.7cm}|>{\raggedright\arraybackslash}m{15.2cm}|}
\caption{Условия включения\hfill\vspace{-0.5\baselineskip}}\label{ksn:tbl1} \\ 
\hhline{|-|-|}
\rowcolor{gray!30}
Режим включения & Условия включения  \\ 
\hhline{|-|-|}
\endfirsthead

\caption*{\hspace{3pt}\emph{Продолжение таблицы \ref{ksn:tbl1}\hfill\vspace{-0.5\baselineskip}}} \\ % сделано по ГОСТ 2.105 п.6.8.7
\hhline{|-|-|}
\rowcolor{gray!30}
Режим включения & Условия включения  \\ 
\hhline{|-|-|}
\endhead
\hhline{|-|-|}

\endfoot
\endlastfoot

ОНЭ1+ННЭ2 &
\makecell[lt]{%
– $U_{\textit{Э1}}$ меньше уставки \textbf{«Uотс Э1»}; \\
– $U_{\textit{Э2}}$ больше уставки \textbf{«Uнал Э2»}. } \\
\hline

ОНЭ2+ННЭ1 &
\makecell[lt]{%
– $U_{\textit{Э1}}$ больше уставки \textbf{«Uнал Э1»}; \\
– $U_{\textit{Э2}}$ меньше уставки \textbf{«Uотс Э2»}. } \\
\hline

КС(УС) &
\makecell[lt]{%
1) Включение c контролем синхронизма: \\
 – $U_{\textit{Э1}}$ больше уставки \textbf{«Uнал Э1»}; \\
 – $U_{\textit{Э2}}$ больше уставки \textbf{«Uнал Э2»}; \\
 – разность напряжений линии и энергообъекта Э1 не превышает по модулю значения \\ заданной уставки \textbf{«dU макс КС(УС)»}; \\
 – разность фаз напряжений линии и энергообъекта Э1 не превышает по модулю \\ значения заданной уставки \textbf{«dФ макс КС»}; \\
 – разность частот напряжений линии и энергообъекта Э1 не превышает значения \\ заданной уставки \textbf{«df макс КС»}; \\
2) Включение с улавливанием синхронизма: \\
 – $U_{\textit{Э1}}$ больше уставки \textbf{«Uнал Э1»}; \\
 – $U_{\textit{Э2}}$ больше уставки \textbf{«Uнал Э2»}; \\
 – разность напряжений линии и энергообъекта Э1 не превышает по \\ модулю значения заданной уставки \textbf{«dU макс КС(УС)»}; \\
 – разность частот напряжений линии и энергообъекта Э1 превышает \\ значение уставки \textbf{«df макс КС»}; \\
 – разность частот напряжений линии и энергообъекта Э1 не превышает значения \\ заданной уставки \textbf{«df пред. доп.»}; \\
 – за время \textbf{«Т опер. вкл.»} с момента подачи сигнала на включение выключателя \\ расхождение углов между векторами напряжений на энергообъектах не превысит $10^\circ$.} \\
\hline

\end{longtable}
\setlength{\extrarowheight}{0.0cm} % возврат параметра высота ячеек в таблице
\normalsize

\item
Если разность частот напряжений шин и линии не превышает уставку \textbf{«df макс КС»}, то разрешается включение с контролем синхронизма. Для режима включения с КС также контролируется:
\begin{itemize}
    \item наличие напряжения на обоих энергообъектах (шины и линия);
    \item разность напряжений шин и линии, не превышающая по модулю значения заданной уставки;
    \item разность фаз напряжений шин и линии, не превышающая по модулю значения заданной уставки.
\end{itemize}

\item
При выводе программного ключа \textbf{«КС»} включение разрешается без контроля частоты, а только по уровню напряжений и разности их фаз.

\item
Если разность частот превышает уставку \textbf{«df макс КС»}, но не выходит за пределы уставки \textbf{«df пред. доп.»}, то разрешается включение с улавливанием синхронизма. Сигнал включения с улавливанием синхронизма формируется с опережением на время включения выключателя. Режим включения с УС можно вывести программным ключом \textbf{«УС»}.

\item
Наличие или отсутствие напряжения на энергообъекте Э1 (линии) контролируется измерительными органами максимального и минимального напряжения, пороги срабатывания которых задаются уставками \textbf{«Uнал Э1»} и \textbf{«Uотс Э1»} соответственно. Если измерительный трансформатор напряжения установлен на линии (определяется положением переключателя \textbf{«Место ТН»}), то контролируется наличие симметричного напряжения. В этом случае дополнительно контролируются уровни напряжений обратной и нулевой последовательности, которые не должны превышать уставок \textbf{«U2 макс КСН»} и \textbf{«3U0 макс КСН»} соответственно.

\item
Пороги срабатывания измерительных органов для контроля напряжения на энергообъекте Э2 (шины) задаются уставками \textbf{«Uнал Э2»} и \textbf{«Uотс Э2»}.

\item
Сигнал разрешения АПВ Э1 \textit{«КСН/ КСН: Разр. АПВ Э1»} формируется при выполнении условий включения в режимах ОНЭ1+ННЭ2, КС или УС. Сигнал разрешения АПВ Э2 \textit{«КСН/ КСН:Разр. АПВ Э2»} формируется при выполнении условий включения в режимах ОНЭ2+ННЭ1, КС или УС. 

\item
Сигнал разрешения оперативного включения \textit{«КСН/ КСН: Разр. опер. вкл»} может формироваться от внутренней логики КСН или от внешнего сигнала в зависимости от положения программного переключателя \textbf{«Контр. опер. вкл.»}.

\item
В положении \textbf{«Внутр.»} оперативное включение разрешается при выполнении условий включения в режимах ОНЭ1+ННЭ2 (при опробовании линии), ОНЭ2+ННЭ1 (при опробовании шин), КС или УС (при замыкании в транзит). Также предусмотрена возможность оперативного включения с контролем отсутствия напряжения на шинах или на линии при введенном программном ключе \textbf{«Опер. вкл. с КОН»}. Оперативное включение может выполняться без контроля напряжений и условий синхронизма при соответствующем заданном режиме АПВ (без контролей) и введенном программном ключе \textbf{«Опер. вкл. без контролей»}.

\item
В положении \textbf{«Внеш.»} оперативное включение разрешается при наличии внешнего сигнала \textit{«КСН/ КСН: Внеш. сигнал КС(УС)»}.

\item
Сигнал разрешения полуавтоматического включения \textit{«КСН/ КСН: Разр. ПАВ»} формируется при выполнении условий включения в режимах КС или УС.

\item
Алгоритм работы КСН приведена на рисунке \ref{pict_ksn:KSN}. Параметры для настройки КСН представлены в таблице \ref{app_set:ksn}.

\medskip
\begin{figure}[H]
\centering
\includegraphics[width=0.85\textwidth,height=0.85\textheight,keepaspectratio]{\fbpath/10.01 АУВ/КСН/img/KSN.pdf}
\captionof{figure}{Функциональная схема логики КСН}\label{pict_ksn:KSN}
\end{figure}


\end{enumerate}